\par
\textbf{Pour chaque affirmation, répondre par VRAI ou FAUX, en justifiant la réponse.}\par
Toute absence de justification ou justification incorrecte ne sera pas prise en compte dans la notation.

Une enquête est réalisée sur l'équipement informatique des 400 élèves d'un collège.
\begin{center}\renewcommand{\arraystretch}{1.35}
\begin{tabular}{|p{3.2cm}|c|c|c|}\hline
& Ont une adresse mail & N'ont pas d'adresse mail & Total\\\hline
Ont un équipement individuel & 50 & 210 & 260\\\hline
N'ont pas d'équipement individuel & 40 & 100 & 140\\\hline
Total & 90 & 310 & 400\\\hline
\end{tabular}
\end{center}

\begin{enumerate}
\item \underline{Affirmation}

50\% des élèves du collège possèdent une adresse mail et un équipement individuel.

\item \underline{Affirmation}

Au moins 50\% des élèves du collège ne possèdent pas d'adresse mail.

\item \underline{Affirmation}

On choisit un élève au hasard dans le collège. La probabilité qu'il ne possède ni adresse mail ni équipement individuel est égale à 25\%.

\item \underline{Affirmation}

Au moins $\dfrac{1}{5}$ des élèves ayant un équipement individuel possèdent également une adresse mail.
\end{enumerate}
